% =========================================================
% Supremum, Infimum, and Bounds
% =========================================================

\section{Supremum and Infimum}\label{sec:sup-inf}

% ---------------------------------------------------------
% TOOLKIT
% ---------------------------------------------------------
\begin{tcolorbox}[colback=gray!6, colframe=gray!40, arc=2pt,
  left=6pt, right=6pt, top=4pt, bottom=4pt,
  title={\small\textbf{Supremum, Infimum, and Bounds — Quick Reference}},
  fonttitle=\small\bfseries]
\small
\begin{tabular}{l l l}
\toprule
\textbf{Concept} & \textbf{Meaning} & \textbf{Detail} \\
\midrule
Upper bound       & $u \geq s$ for all $s \in S$                       & \hyperref[def:bounds]{↓ Def} \\
Lower bound       & $\ell \leq s$ for all $s \in S$                    & \hyperref[def:bounds]{↓ Def} \\
Supremum          & Least upper bound of $S$                           & \hyperref[def:sup-inf]{↓ Def} \\
Infimum           & Greatest lower bound of $S$                        & \hyperref[def:sup-inf]{↓ Def} \\
Uniqueness        & Suprema and infima are unique when they exist       & \hyperref[prop:sup-unique]{↓ Prop} \\
Max/min vs.\ sup/inf & Maximum $\Rightarrow$ supremum; not conversely & \hyperref[rem:max-vs-sup]{↓ Rem} \\
Duality           & $\sup$ in $(A,\leq)$ = $\inf$ in $(A,\geq)$        & \hyperref[prop:sup-inf-duality]{↓ Prop} \\
\bottomrule
\end{tabular}
\end{tcolorbox}

\vspace{1em}

% ---------------------------------------------------------
% Motivation
% ---------------------------------------------------------

Every student of analysis encounters the least upper bound property of
$\mathbb{R}$ early and prominently.  But this property is not a peculiarity
of real numbers: it is an order-theoretic concept that can be formulated in
any partially ordered set.  Before we reach $\mathbb{R}$, it is essential to
understand supremum and infimum in their natural habitat — abstract ordered
sets — so that when the completeness axiom appears in Volume~II, you will
recognise exactly what it is asserting and why $\mathbb{Q}$ fails to satisfy it.

The definitions in this section are also the backbone of measure theory.
The extended real line $\overline{\mathbb{R}} = \mathbb{R} \cup \{-\infty,
+\infty\}$, equipped with the usual order, is the ambient poset in which
measures take their values.  Countable additivity, the monotone convergence
theorem, and Fatou's lemma are all statements about suprema of sequences.
Understanding supremum and infimum at the abstract level here means those
later theorems will feel inevitable rather than arbitrary.

% ---------------------------------------------------------
% Upper and Lower Bounds (restate cleanly here for self-containment)
% ---------------------------------------------------------

The definition of upper and lower bound was given in
\hyperref[def:bounds]{the Ordered Sets section}.  For convenience
we recall: $u \in A$ is an \emph{upper bound} of $S \subseteq A$ if
$\forall s \in S,\ s \leq u$; a \emph{lower bound} satisfies the
reverse.  We say $S$ is \emph{bounded above} (resp.\ \emph{bounded below},
\emph{bounded}) if it has at least one upper (resp.\ lower, both) bound.

\begin{remark}[Bounds depend on the ambient poset]
\label{rem:bounds-ambient}
Whether $S$ is bounded above depends on the poset $(A, \leq)$, not on
$S$ alone.  The set $\{r \in \mathbb{Q} : r^2 < 2\}$ is bounded above in
$(\mathbb{Q}, \leq)$ — the rational number $2$ is an upper bound — but
it has \emph{no} supremum inside $\mathbb{Q}$.  The same set in
$(\mathbb{R}, \leq)$ has supremum $\sqrt{2} \in \mathbb{R}$.  Being
bounded above is necessary but not sufficient for a supremum to exist;
the completeness axiom for $\mathbb{R}$ is precisely the assertion that
sufficiency holds in $\mathbb{R}$.
\end{remark}

% ---------------------------------------------------------
% Supremum and Infimum
% ---------------------------------------------------------

\begin{tcolorbox}[colback=propbox, colframe=propborder, arc=2pt,
  left=6pt, right=6pt, top=4pt, bottom=4pt,
  title={\small\textbf{Definition (Supremum and Infimum)}},
  fonttitle=\small\bfseries]
\label{def:sup-inf}
Let $(A, \leq)$ be a poset and $S \subseteq A$.

The \emph{supremum} (or \emph{least upper bound}) of $S$, written
$\sup S$, is an element $u^* \in A$ satisfying:
\begin{enumerate}[label=(\roman*)]
  \item $u^*$ is an upper bound of $S$: $\forall s \in S,\ s \leq u^*$;
  \item $u^*$ is the \emph{least} upper bound: if $u$ is \emph{any}
        upper bound of $S$, then $u^* \leq u$.
\end{enumerate}

The \emph{infimum} (or \emph{greatest lower bound}) of $S$, written
$\inf S$, is an element $\ell^* \in A$ satisfying:
\begin{enumerate}[label=(\roman*)]
  \item $\ell^*$ is a lower bound of $S$: $\forall s \in S,\ \ell^* \leq s$;
  \item $\ell^*$ is the \emph{greatest} lower bound: if $\ell$ is \emph{any}
        lower bound of $S$, then $\ell \leq \ell^*$.
\end{enumerate}
\end{tcolorbox}

\begin{remark*}[Fully quantified form]
$u^* = \sup S \;\leftrightarrow\; (\forall s \in S,\; s \leq u^*) \;\wedge\;
(\forall u \in A,\; (\forall s \in S,\; s \leq u) \;\rightarrow\; u^* \leq u)$.

$\ell^* = \inf S \;\leftrightarrow\; (\forall s \in S,\; \ell^* \leq s) \;\wedge\;
(\forall \ell \in A,\; (\forall s \in S,\; \ell \leq s) \;\rightarrow\; \ell \leq \ell^*)$.
\end{remark*}

\begin{remark}[How to read these definitions]
\label{rem:sup-reading}
Condition (i) says $u^*$ is a \emph{candidate}: it stands above everything in
$S$. Condition (ii) says $u^*$ is the \emph{best} candidate: it is at most as
large as any other upper bound.  Together they say: $u^*$ is the unique
element of $A$ that is simultaneously an upper bound of $S$ and no larger than
any other upper bound.

A useful restatement of condition (ii) is its contrapositive: if $v < u^*$,
then $v$ is \emph{not} an upper bound of $S$, meaning there exists some
$s \in S$ with $v < s$.  This is how the condition is typically used in
proofs in analysis: to verify that $u^* = \sup S$, you show (i) $u^*$
is an upper bound, and (ii) for every $\varepsilon > 0$, the value
$u^* - \varepsilon$ is \emph{not} an upper bound, i.e., there exists
$s \in S$ with $s > u^* - \varepsilon$.
\end{remark}

% ---------------------------------------------------------
% Uniqueness
% ---------------------------------------------------------

\begin{proposition}[\texorpdfstring{\hyperref[prf:sup-unique]{Uniqueness of Supremum and Infimum}}{Uniqueness of Supremum and Infimum}]
\label{prop:sup-unique}
Let $(A, \leq)$ be a poset and $S \subseteq A$.  If $\sup S$ exists, it is
unique.  If $\inf S$ exists, it is unique.
\end{proposition}

\begin{remark*}[Fully quantified form]
$\forall u^*, v^* \in A$: if both $u^*$ and $v^*$ satisfy the two conditions
in the definition of $\sup S$, then $u^* = v^*$.
\end{remark*}

\begin{remark}[Proof idea]
If $u^*$ and $v^*$ are both suprema, then $u^*$ is an upper bound, so
the minimality of $v^*$ gives $v^* \leq u^*$; and symmetrically $u^* \leq v^*$.
Antisymmetry of the partial order then forces $u^* = v^*$.  This is a
standard antisymmetry argument: when two elements are mutually $\leq$-related
in a partial order, they must be equal.
\end{remark}

% ---------------------------------------------------------
% Sup need not lie in S; max vs sup
% ---------------------------------------------------------

\begin{remark}[Supremum need not belong to $S$]
\label{rem:max-vs-sup}
The set $S = (0,1) \subseteq \mathbb{R}$ is open: it contains neither
endpoint.  Yet $\sup S = 1$ and $\inf S = 0$.  Neither $0$ nor $1$
belongs to $S$.

When the supremum does happen to lie in $S$, it is called the
\emph{maximum} of $S$, written $\max S$.  Thus $\max S = \sup S$ when
$\sup S \in S$, and $\max S$ is undefined when $\sup S \notin S$.
Every maximum is a supremum, but the converse fails: a supremum may
exist without being a maximum.  The same relationship holds between
infimum and minimum.

This distinction is critical in analysis.  A continuous function on a
closed bounded interval always achieves its maximum (the Extreme Value
Theorem), but on an open interval it may approach a supremum without
ever attaining it.
\end{remark}

% ---------------------------------------------------------
% Sup need not exist
% ---------------------------------------------------------

\begin{remark}[Supremum need not exist]
\label{rem:sup-nonexistence}
Being bounded above is necessary but not sufficient for a supremum to exist.
Consider $S = \{r \in \mathbb{Q} : r^2 < 2\}$ in the poset $(\mathbb{Q},\leq)$.
Every element of $S$ satisfies $r < 2$, so $2 \in \mathbb{Q}$ is an upper
bound.  But the candidate supremum would need to be $\sqrt{2}$, and
$\sqrt{2} \notin \mathbb{Q}$.  Any rational upper bound $q > \sqrt{2}$ is not
the least upper bound, because $\frac{q + \sqrt{2}}{2}$ is a smaller upper
bound (and is also irrational, which means we can always find a rational
strictly between any rational upper bound and $\sqrt{2}$, and that rational is
still an upper bound).  In short: $S$ is bounded above in $\mathbb{Q}$, but
$\sup S$ does not exist in $\mathbb{Q}$.

This is the precise failure that the \emph{completeness axiom} for
$\mathbb{R}$ corrects: it asserts that every nonempty subset of $\mathbb{R}$
that is bounded above has a supremum in $\mathbb{R}$.  The rational numbers
satisfy every other ordered field axiom but fail this one.
\end{remark}

% ---------------------------------------------------------
% Examples
% ---------------------------------------------------------

\begin{example}[Sup and inf in divisibility]
\label{ex:sup-divisibility}
Order the set $D = \{1,2,3,4,6,12\}$ by divisibility: $a \leq b$ iff $a \mid b$.
Take $S = \{4, 6\}$.

\emph{Upper bounds of $S$}: an element $u \in D$ must satisfy $4 \mid u$ and
$6 \mid u$, i.e., $u$ must be a common multiple of $4$ and $6$ in $D$.
The only such element is $12$, so $\sup S = 12 = \mathrm{lcm}(4,6)$.

\emph{Lower bounds of $S$}: an element $\ell \in D$ must satisfy $\ell \mid 4$
and $\ell \mid 6$, i.e., $\ell$ must be a common divisor of $4$ and $6$ in $D$.
The common divisors are $\{1,2\}$; the greatest is $2$, so $\inf S = 2 =
\gcd(4,6)$.

This example reveals a deep pattern: in any divisibility poset, $\sup$ computes
the least common multiple and $\inf$ computes the greatest common divisor.
Order theory and number theory are speaking the same language.
\end{example}

\begin{example}[Sup and inf for function spaces]
\label{ex:sup-functions}
Let $B(X)$ be the set of bounded real-valued functions on a set $X$, ordered
pointwise: $f \leq g$ iff $f(x) \leq g(x)$ for all $x \in X$.  For a
collection $\{f_\alpha\}$ of bounded functions, the pointwise supremum is
$(\sup_\alpha f_\alpha)(x) = \sup_\alpha f_\alpha(x)$.  In measure theory and
integration theory, one constantly takes suprema and infima of sequences of
functions; this example shows that such operations are instances of the
abstract sup/inf defined here.
\end{example}

% ---------------------------------------------------------
% Useful characterisation
% ---------------------------------------------------------

\begin{proposition}[\texorpdfstring{\hyperref[prf:sup-char]{Characterisation of Supremum}}{Characterisation of Supremum}]
\label{prop:sup-char}
Let $(A, \leq)$ be a poset, $S \subseteq A$, and $u^* \in A$.  Then
$u^* = \sup S$ if and only if:
\begin{enumerate}[label=(\roman*)]
  \item $u^*$ is an upper bound of $S$; and
  \item for every $v \in A$ with $v < u^*$, there exists $s \in S$ with
        $v < s$.
\end{enumerate}
\end{proposition}

\begin{remark*}[Fully quantified form]
$u^* = \sup S \;\leftrightarrow\; (\forall s \in S,\; s \leq u^*) \;\wedge\;
(\forall v \in A,\; v < u^* \;\rightarrow\; \exists s \in S,\; v < s)$.
\end{remark*}

\begin{remark}[Why this characterisation is useful]
\label{rem:sup-char-utility}
Condition (ii) here is the contrapositive of the minimality condition in
the definition.  In analysis this is the working form: to check that
$u^*$ is the supremum, you verify that nothing strictly below $u^*$ is an
upper bound.  In $\mathbb{R}$ this becomes: for every $\varepsilon > 0$,
there exists $s \in S$ with $s > u^* - \varepsilon$.  This is the exact
form used in the proof that $\sup$ interacts correctly with limits,
integrals, and measures.
\end{remark}

% ---------------------------------------------------------
% Duality
% ---------------------------------------------------------

\begin{proposition}[\texorpdfstring{\hyperref[prf:sup-inf-duality]{Duality of Supremum and Infimum}}{Duality of Supremum and Infimum}]
\label{prop:sup-inf-duality}
Let $(A, \leq)$ be a poset and $S \subseteq A$.  Then
$\inf_{(A,\leq)} S = \sup_{(A,\geq)} S$,
where $(A,\geq)$ denotes the dual poset.

In particular: if every nonempty bounded-above subset of $(A,\leq)$ has a
supremum, then every nonempty bounded-below subset of $(A,\leq)$ has an
infimum.
\end{proposition}

\begin{remark*}[Fully quantified form]
$\forall S \subseteq A$: $\ell^*$ is the infimum of $S$ in $(A,\leq)$
$\;\leftrightarrow\;$ $\ell^*$ is the supremum of $S$ in $(A,\geq)$.
\end{remark*}

\begin{remark}[Consequence for $\mathbb{R}$]
\label{rem:inf-from-completeness}
This proposition explains why the completeness axiom for $\mathbb{R}$ is
stated only for suprema: the infimum statement comes for free by duality.
The axiom says every nonempty bounded-above subset has a supremum; the
proposition immediately gives that every nonempty bounded-below subset has
an infimum.  You will use this constantly without comment in Volume~II.
\end{remark}

% ---------------------------------------------------------
% Connection to sup/inf of sequences
% ---------------------------------------------------------

\begin{remark}[Sequences and sup/inf]
\label{rem:seq-sup-inf}
A sequence $(a_n)_{n=1}^{\infty}$ in $\mathbb{R}$ is a function
$\mathbb{N} \to \mathbb{R}$, so its range $\{a_n : n \in \mathbb{N}\}$ is a
subset of $\mathbb{R}$.  The quantities $\sup_n a_n$ and $\inf_n a_n$ are
the supremum and infimum of this range set in the poset $(\mathbb{R}, \leq)$.

More subtle are the \emph{limit superior} and \emph{limit inferior}:
\[
  \limsup_{n\to\infty} a_n = \inf_{n \geq 1} \sup_{k \geq n} a_k,
  \qquad
  \liminf_{n\to\infty} a_n = \sup_{n \geq 1} \inf_{k \geq n} a_k.
\]
These are iterated suprema and infima.  The fact that they always exist in
$\overline{\mathbb{R}}$ (the extended reals, where $\pm\infty$ are adjoined
to supply missing bounds) is a direct application of the completeness of
$\mathbb{R}$.  Mastering the abstract definitions here makes those formulas
transparent rather than mysterious.
\end{remark}

\begin{remark}[Transition]
\label{rem:transition-to-lattice}
Supremum and infimum are defined for arbitrary subsets of a poset, but they
need not always exist.  When \emph{every} two-element subset $\{a,b\}$ has
both a supremum and an infimum, the poset has enough structure to define
binary operations $a \vee b = \sup\{a,b\}$ and $a \wedge b = \inf\{a,b\}$.
This additional structure is called a \emph{lattice}, and it is developed in
the next section.  Lattices are the precise abstract setting for the
$\sigma$-algebras and topologies encountered in measure theory and topology.
\end{remark}
